\begin{table}[htbp]
\centering
\caption{Rendimento da blocagem nos três passos exportados do OpenRecLink}
\label{tab:blocagem-rendimento}
\small
\setlength{\tabcolsep}{4pt}
\begin{tabular}{crrrrr}
\toprule
Passo & \shortstack{Potenciais\\pares} & Vínculos & \shortstack{Proporção de\\vínculos (\%)} & \shortstack{Mediana da\\nota final} & \shortstack{Vínculos na\\área cinza (\%)} \\
\midrule
6  & 2.649  & 80  & 3,0  & 6,48 & 22,5 \\
9  & 1.449  & 127 & 8,8  & 4,64 & 82,7 \\
11 & 57.598 & 40  & 0,07 & 4,42 & 70,0 \\
\bottomrule
\end{tabular}

\smallskip
{\footnotesize\textit{Fonte}: elaboração própria, auditoria M8 do extrato do OpenRecLink.
\par\textit{Nota}: potenciais pares e vínculos são contados distintamente em cada passo; a proporção corresponde ao número de vínculos dividido pelo número de potenciais pares do passo. A mediana da nota final é calculada sobre todas as linhas de cada passo. A área cinza da nota corresponde à nota final do Comparador de Registros (CdR) em $[5,9)$, com a fração desta tabela calculada entre os vínculos do passo; não corresponde à região abaixo do corte do OpenRecLink.}
\end{table}
